% Generated draft: requires mathematical and source review.
\documentclass[11pt]{article}
\usepackage[margin=25mm]{geometry}
\usepackage{fontspec}
\setmainfont{Latin Modern Roman}
\usepackage{amsmath,amssymb,mathtools,unicode-math}
\setmathfont{Latin Modern Math}
\usepackage{enumitem,needspace,xurl,hyperref}
\setlength{\parindent}{0pt}
\setlength{\parskip}{5pt}
\setlength{\emergencystretch}{3em}
\newcommand{\N}{\mathbb{N}}
\newcommand{\Z}{\mathbb{Z}}
\newcommand{\Q}{\mathbb{Q}}
\newcommand{\R}{\mathbb{R}}
\newcommand{\C}{\mathbb{C}}
\newcommand{\F}{\mathbb{F}}
\newcommand{\A}{\mathbb{A}}
\newcommand{\PP}{\mathbb{P}}
\newcommand{\E}{\mathbb{E}}
\newcommand{\eps}{\varepsilon}
\newcommand{\dd}{\,\mathrm{d}}
\begin{document}
\section*{UnsolvedMath ID: 30001643}
\textbf{Source problem code:} catalogue entry
\subsection{Definitions and mathematical statement}
Let \(H\) be a Hilbert space and let \(L_H:D(L_H)\subseteq H\to H\) denote the operator called \(L\) in the supplied record. Write \(\sigma_H(L_H)\) for its spectrum and \(\ker L_H=\{f\in D(L_H):L_Hf=0\}\). Let \(T_H(t)=e^{tL_H}\) denote the semigroup appearing in the source assertion, and let \(\Pi_H\) denote the operator written there as \(\Pi\). The record does not explicitly define this operator or specify the operator-domain and semigroup-generation hypotheses.

The supplied Hilbert-space assumptions are that there exist \(a<0\) and a finite constant \(C_a>0\) such that
\[
\sigma_H(L_H)\cap\{z\in\C:\operatorname{Re}z>a\}=\{0\},
\qquad \dim\ker L_H<\infty,
\]
and, with the usual forward-time interpretation of the decay assertion,
\[
\|T_H(t)f-\Pi_Hf\|_H\leq C_a e^{at}\|f\|_H
\qquad(f\in H,\ t\geq0).
\]
The source does not specify whether the scalar field is complex or whether complexification is intended when defining the spectrum of a real operator. The role of \(\Pi_H\) as an equilibrium projection, including its range and whether it is a spectral or orthogonal projection, requires clarification from the source.

Let \(E\) be a Banach space containing \(H\), and let \(L_E:D(L_E)\subseteq E\to E\) extend \(L_H\). Under the identification \(H\subseteq E\), operator extension means
\[
D(L_H)\subseteq D(L_E),\qquad L_Ef=L_Hf\quad(f\in D(L_H)).
\]
No continuity or density assumption on the inclusion \(H\subseteq E\) is supplied.

\textbf{Extension question.} Determine additional hypotheses on the inclusion, the operators and their semigroups under which the Hilbert-space spectral and decay conclusions transfer to \(E\). A literal same-rate formulation of the requested conclusions is
\[
\sigma_E(L_E)\cap\{z\in\C:\operatorname{Re}z>a\}=\{0\},
\qquad \dim\ker L_E<\infty,
\]
together with a semigroup \(T_E(t)=e^{tL_E}\), an appropriate bounded equilibrium projection \(\Pi_E\), and a finite constant \(C_{a,E}>0\) satisfying
\[
\|T_E(t)f-\Pi_Ef\|_E\leq C_{a,E}e^{at}\|f\|_E
\qquad(f\in E,\ t\geq0).
\]
The intended compatibility conditions would need to specify whether
\[
T_E(t)|_H=T_H(t),\qquad \Pi_E|_H=\Pi_H,
\qquad \operatorname{Ran}\Pi_E=\ker L_E
\]
are required. These conditions express a possible precise transfer target; they are not additional hypotheses explicitly supplied by the record. Likewise, the record does not settle whether the same exponent \(a\) is required or whether transfer with a different negative exponent is sufficient. The question asks for sufficient hypotheses, rather than asserting transfer for every extension. The supplied verification note says that the report gives sufficient splitting conditions and applications, but those conditions are not included in the supplied record and cannot be reconstructed here.
\subsection{Short English statement}
The problem asks when spectral isolation of zero and exponential convergence to equilibrium for an operator on a Hilbert space persist after extending the operator to a larger Banach space. The embedding, generator assumptions, equilibrium projection and intended decay rate need clarification. The supplied curation labels the problem partially solved; this draft does not independently verify that assessment.
\subsection{Sources}
\begin{itemize}
\item Local project registry: supplied complete source record for permanent UnsolvedMath ID \(30001643\), source problem code \texttt{catalogue entry}, title \emph{Transferring Hypocoercivity to Larger Banach Spaces}, in the Oberwolfach Reports Open Problems collection. The record gives proposed year 2010 and names no proposer.
\item Supplied report citation: \emph{Classical and Quantum Mechanical Models of Many-Particle Systems} (2011), Oberwolfach Reports, pp.\ 3159--3236, DOI \texttt{10.4171/owr/2010/54}. Supplied source URL and literature-source label ``original Oberwolfach problem report'': \url{https://doi.org/10.4171/owr/2010/54}.
\item Supplied literature source labeled ``preserving hypocoercivity on larger spaces'': \url{https://www.ceremade.dauphine.fr/~mischler/expo/oberwolfach2010.pdf}. No author or fuller bibliographic citation is supplied for this item.
\item Literature assessment dated exactly \texttt{2026-08-21}: the supplied curation labels the problem partially solved and reports that hypocoercivity and spectral transfer have been proved for several larger weighted Banach spaces, while a general extension question for arbitrary \(E\) and \(L\) remains open. This is supplied dated curation, not an independently verified literature assessment.
\item The supplied statement-status label is \texttt{corrected\_verified}. Its accompanying verification note attributes the Hilbert-space assumptions, extension question, sufficient splitting conditions and applications to the official report. This attribution is preserved as supplied; the linked sources have not been independently consulted for this draft.
\end{itemize}
\end{document}
